\documentclass[12pt]{article}
\usepackage[utf8]{inputenc}
\usepackage[T1]{fontenc}
\usepackage{hyperref}
\title{Urumi Language Resources for K-12 Curriculum Development}
\author{Compiled by Claude Sonnet 5 (claude-sonnet-5) for Andr\'as Kornai}
\date{}
\begin{document}
\maketitle

\section{Introduction}
Urumi (ISO 639-3 uru) is an extinct Tupian language of Rond\^onia, Brazil, placed in the Ramarama branch. Sources disagree on when it died out, giving either after 1913 or the latter half of the 20th century. The record consists of wordlists and numerals, including the Rondon Commission glossary.

\section{Primary Educational Resources}
None exist, as far as search results show.

\section{Linguistic and Reference Materials}
\begin{itemize}
\item Wikipedia entry, which includes the numerals: \url{https://en.wikipedia.org/wiki/Urumi_language}
\item Ramarama languages: \url{https://en.wikipedia.org/wiki/Ramarama_languages}
\item Numerals database entry for Urumi: \url{https://lingweb.eva.mpg.de/channumerals/Urumi.htm}
\end{itemize}
No grammar or dictionary was found.

\section{Religious and Community Resources}
None found.

\section{Digital and Audio Resources}
None found. No recordings are known.

\section{Curriculum Development Context}
The documentation is too thin for language revitalization teaching. A heritage curriculum could focus on the history of the people, using Portuguese as the language of instruction, with the attested numerals and words as a small cultural component. Any such work needs descendant communities and Brazilian indigenous education authorities.
\end{document}